% Imporditud: tools/import_eksperimendid.py
% Allikas: Eesti füüsikaolümpiaadi eksperimendi ülesannete
%   kogu (Rael Kalda, 2023), ülesanne Ü36, lahendus L36.
\setAuthor{EFO žürii}
\setRound{piirkonnavoor}
\setYear{1998}
\setNumber{P E1}
\setDifficulty{2}
\setTopic{varia}
\setVahendid{leht paberit, mõõtejoonlaud}
\setPunktid{10}
\setAllikas{Eksperimendi kogu Ü36, lahendus lk 49}
\setLahendusseis{toores}

\prob{Paberileht}
Paberi pakendil on kiri ``80 g/m$^2$''. Määrata antud paberilehe mass ja paksus. Kirjeldada ülesande lahendust.
\textit{Märkus:} Sama ülesanne oli aastal 2012 põhikooli piirkonnavoorus.

\hint

\solu
Teooria: Kiri ``80 g/m$^2$'' tähendab, et paberi iga ruutmeetri mass on 80 g (1 p). 80 g/m$^2$ võib tähistada, näiteks, tähega r. Kasutades ühikut, saab tuletada valemi r = m/S (1 p), kus m on paberilehe mass ja S on pindala. Paberilehe pindala S = ab, kus a ja b on paberilehe mõõtmed. Paberilehe mass m = rS (1 p). Paberilehe paksuse määramiseks voltida leht kokku ja mõõta joonlauaga saadud paki paksus, paki paksus jagada paberikihtide arvuga (1 p). Paberilehe paksus d = L/n, kus L on paki paksus ja n on paberi kihtide arv pakis (1 p).

Mõõtmised: Paberilehe pikkuse ja laiuse mõõtmine (1 p). Paberipaki paksuse mõõtmine (1 p).

Arvutused: Paberilehe pindala (1 p). Paberilehe mass koos ühikuga (1 p). Paberilehe paksus koos ühikuga (1 p).
\probend
